\section{Separate accepted audio from published metadata}
Maintain flags for whether format and data have been seen, accepted format fields,
the candidate data offset and its length, and the next chunk offset. Finding data
proves a possible audio region; it does not finish the outer walk. Requiring the format
first lets the parser check frame divisibility at that point. Continue validating later
chunks before copying the final descriptor to caller memory.

The worked input has sixteen PCM bytes and four frames. Deleting its last pad does
not change those sample bytes. If the outer size remains 74, the file already fails
its exact extent; if the size is adjusted to 73, its final odd-sized chunk instead
lacks the required pad. Neither failure may leave partially accepted rate or channel fields in the
caller descriptor. Use an untouched sentinel descriptor to observe that distinction.

After a successful parse, the program can form the PCM pointer from the particular
buffer base plus 56 and pass four frames through its description prefix. There is no
sample decompression step for the supported PCM representation. The processing kernel
still has to interpret signed little-endian sample codes and interleaving correctly.
Validation authorizes that slice; it does not silently transform it into planar
channels or a different byte order.
